% Headers file

% Standard packages
\usepackage[T1]{fontenc}
\usepackage{graphicx}
\usepackage{color}
\usepackage{epstopdf}               % Auto convert eps to pdf for pdflatex
\usepackage{amsmath}                % Support math
\usepackage{booktabs}               % Required by siunitx, adds toprule, midrule, and bottomrule lines
%\usepackage{cancel}                 % Required by siunitx
\usepackage{caption}                % Required by siunitx
%\usepackage{cleveref}               % Required by siunitx
%\usepackage{colortbl}               % Required by siunitx
%\usepackage{helvet}                 % Required by siunitx
%\usepackage{mathpazo}               % Required by siunitx (changes font to libertine)
\usepackage{pgfplots}               % Required by siunitx
\usetikzlibrary{shapes.geometric}   % Diamond decision nodes (generated flow charts)
\usetikzlibrary{arrows.meta}        % Stealth arrow tips (generated diagrams)
\usetikzlibrary{positioning}        % Relative node placement (generated diagrams)
\usepackage{xcolor}                 % Required by siunitx
\usepackage{listings}               % Code listings
\usepackage{siunitx}                % SI unit formatting
%\usepackage{upgreek}                % Better greek symbols (not needed with libertinust1math)
%\sisetup{math-micro=\mu}            % Uses a better mu symbol for siunitx (not needed with libertinust1math)
\usepackage{rotating}               % Rotate figures
\usepackage{longtable}              % Split table between multiple pages
\usepackage{mathrsfs}               % Fancy cursive math characters
\usepackage{subcaption}             % Newest/best way to add multiple figures in one figure
\usepackage{cite}                   % Fancy ordering of multi-citations
\usepackage{colortbl}
%\usepackage{float}                  % Figures that can float around the document
\usepackage{placeins}               % \FloatBarrier -- flush pending floats at a boundary
\usepackage{svg}
%\usepackage[off]{svg-extract}
\usepackage{relsize}

% Extra packages
\usepackage{acronym}                % Glossary of terms and abbreviations (chapter 1)
\usepackage{titlesec}               % Compact chapter and part headings (report class)
\usepackage{tcolorbox}              % Note and Caution callout boxes
\usepackage{verbatim}
\usepackage{fancyhdr}
\pagestyle{fancy}
\usepackage{lastpage}
\usepackage{amssymb}
\usepackage{gensymb}
\usepackage{enumitem}
\usepackage{mathtools}
\usepackage{bm}
\usepackage{hhline}
\usepackage{multirow}
\usepackage{multicol}               % Two column code listings

\usepackage{tikz}                   % Graphic element creator (needed for circuitikz)
\usetikzlibrary{calc}
%\usepackage{tikzit}	% Enables support of the tikzit GUI for making tikz figures. Requires the tikzit.sty file to be in the root latex directory
\usepackage[american, siunitx, arrowmos, smartlabels, betterproportions, emptypmoscircle]{circuitikz} % Makes circuit schematics
\usepackage{tikz-timing}
\usetikztiminglibrary[new={char=Q,reset char=R}]{counters}
\usetikztiminglibrary{nicetabs} % a bit strange with \Huge; use belowrulesep to adjust

\usepackage{amsmath}

\usepackage{ltablex}
\usepackage[margins=raggedright]{floatrow}

\usepackage{bytefield}
\usepackage{soul}	% For highlighting with \hi command
\usepackage{textcomp}

% Fonts
\usepackage{libertine}	% Libertine (the main font)
\usepackage[lcgreekalpha]{libertinust1math}	% Libertine (as the math font)
\usepackage[varl]{inconsolata}	% A monospace font similar to Consolas. The varl option makes the lowercase l look different than a numeral 1.
%\usepackage{cascadia-code}

% Last packages
\usepackage{csquotes}               % Required by siunitx
\usepackage{hyperref}               % Hyperlinks; the link colours are set after the theme colours below


% declare the path(s) where your graphic files are
\graphicspath{{./}, {./figures/}}
% and their extensions so you won't have to specify these with
% every instance of \includegraphics
\DeclareGraphicsExtensions{.pdf,.eps,.svg,.SVG,.png,.PNG,.jpeg,.jpg,.JPEG,.JPG}

% ---------------------------------------------------------------------------
% FIGURE THEME.
% Every generated block diagram in this manual draws from this one palette and
% this one style set, so the figures read as one hand-made set rather than as
% thirty independent drawings.
% The theme is the whole-chip figure's (Figure 1): a red package boundary and
% red accents, greyscale blocks, sans-serif labels throughout.
% vestaRed is exactly red!75!black and vestaRedText exactly red!70!black, which
% is what that figure was already drawn in.
% Emitters should reach for a v-prefixed style below rather than spelling a
% fill or a line width inline. A one-off that genuinely needs its own value
% should still be written as a modification of one of these.
% ---------------------------------------------------------------------------
\definecolor{vestaRed}{rgb}{0.75,0,0}
\definecolor{vestaRedText}{rgb}{0.70,0,0}
\definecolor{vestaInk}{rgb}{0,0,0}
\tikzset{
	% Boxes. Three weights of fill, one stroke, one corner radius.
	vblock/.style={thick, rounded corners=2pt, draw=vestaInk, fill=black!8},
	vblocklt/.style={thick, rounded corners=2pt, draw=vestaInk, fill=black!3},
	vblockw/.style={thick, rounded corners=2pt, draw=vestaInk, fill=white},
	vblockem/.style={line width=1.1pt, rounded corners=2pt, draw=vestaInk, fill=black!12},
	% The bus bar, and the package boundary.
	vbar/.style={thick, draw=vestaInk, fill=black!15},
	vbound/.style={draw=vestaRed, line width=1.2pt},
	% A grouping region. Never a full-bleed grey band: a thin dashed outline
	% that leaves the paper white, because a solid band behind a figure is the
	% single thing that made these drawings look machine-made.
	vregion/.style={draw=black!35, line width=0.5pt, dashed, rounded corners=3pt},
	% Text. Sans throughout; italic is for asides only, never for a block name.
	vhd/.style={font=\sffamily\scriptsize, align=center, inner sep=1pt, anchor=north},
	vbc/.style={font=\sffamily\scriptsize, align=center, inner sep=1pt},
	vsm/.style={font=\sffamily\tiny, align=center, inner sep=1pt},
	vnote/.style={font=\sffamily\scriptsize\itshape, text=black!55, align=left, inner sep=1.5pt},
	vgroup/.style={font=\sffamily\small\itshape, text=black!55, align=left, fill=white, inner sep=1.5pt},
	vredlab/.style={font=\sffamily\small\bfseries, text=vestaRedText, align=left},
	% Lines. One arrow head for the whole manual.
	vwire/.style={semithick, draw=vestaInk},
	vflow/.style={->, >=Stealth, semithick, draw=vestaInk},
	vbus/.style={<->, >=Stealth, thick, draw=vestaInk},
	vaccent/.style={->, >=Stealth, semithick, draw=vestaRed},
	vghost/.style={dashed, draw=black!45, semithick},
}

% Cross-references and register names are visibly clickable, in the theme's
% text red. Set here rather than in the \usepackage line because the colour is
% defined above, after hyperref is loaded.
\hypersetup{colorlinks=true, linkcolor=vestaRedText, citecolor=vestaRedText, urlcolor=vestaRedText, filecolor=vestaRedText}

% Running head needs 14pt with the Libertine face; fancyhdr warns on every page otherwise.
\setlength{\headheight}{14pt}

% Callouts. Exactly two tiers: Note (advisory) and Caution (the hardware
% misbehaves if the text is ignored). Both draw from the theme colours only.
\tcbset{
	vestanote/.style={colback=black!3, colframe=black!35, coltitle=black, fonttitle=\sffamily\bfseries, boxrule=0.5pt, arc=2pt, left=4pt, right=4pt, top=2pt, bottom=2pt, before skip=6pt, after skip=8pt},
	vestacaution/.style={colback=black!3, colframe=vestaRed, coltitle=vestaRedText, fonttitle=\sffamily\bfseries, boxrule=0.8pt, arc=2pt, left=4pt, right=4pt, top=2pt, bottom=2pt, before skip=6pt, after skip=8pt},
}
\newtcolorbox{Note}{vestanote, title=Note}
\newtcolorbox{Caution}{vestacaution, title=Caution}

% Heading-level shift for an included file that still uses article levels.
% Use inside a group: {\ShiftHeadingsUp \input{file}}. The per-implementation
% analog chapter is the one consumer; it is not owned by this template.
\newcommand{\ShiftHeadingsUp}{\let\subsubsection\subsection \let\subsection\section \let\section\chapter}

% Table highlight color
\definecolor{tablehighlightcolor}{gray}{0.925}
\definecolor{linuxhighlightcolor}{gray}{0.85}
\sethlcolor{linuxhighlightcolor}

% A handy macro for the address space table
\newcommand{\memsection}[4]{%
	% define the height of the memsection
	\bytefieldsetup{bitheight=#3\baselineskip}%
	\bitbox[]{10}{%
		\texttt{#1} % print start address
		\\
		%   do some spacing
		\vspace{#3\baselineskip}
		\vspace{-2\baselineskip}
		\vspace{-#3pt}
		\texttt{#2} % print end address
	}%
	\bitbox{16}{\sffamily\small #4}	% print box with caption, sans (figure theme)
}

% Add a macro for tikz-timing time axis (from https://tex.stackexchange.com/questions/67821/tikz-timing-is-there-a-straightforward-way-to-add-a-numbered-time-axis)
\newcommand{\timingaxis}[2][]{% arg 1 is a default argument (not given), argument #2 is the time increment. Call with \timingaxis{time_increment}
	\begin{scope}[#1]
		\draw [timing/table/axis] (0,-\nrows-1) -- (\twidth+1,-\nrows-1);
		\foreach \n in {0,1,...,\twidth} {
			\draw [timing/table/axis ticks] (\n,-\nrows-1+.1) -- +(0,-.2);
		}
		\foreach \n in {0,#2,...,\twidth} {
			\draw (\n,-\nrows-1-.1) node [below,inner sep=2pt] {\scriptsize\n};
		}
	\end{scope}
}
\tikzset{%
	timing/table/axis/.style={->,>=latex},
	timing/table/axis ticks/.style={},
}

% Listing definitions
\definecolor{lstgreen}{rgb}{0,0.6,0}
\definecolor{lstgray}{rgb}{0.4,0.4,0.4}
\lstdefinestyle{customC}{
	language=C,
	basicstyle=\scriptsize\ttfamily,
	numbers=left,
	numberstyle=\footnotesize,
	keywordstyle=\color{blue},
	identifierstyle=\color{black},
	commentstyle=\color{lstgreen},
	stringstyle=\color{lstgray},
	stepnumber=1,
	numbersep=10pt,
	backgroundcolor=\color{white},
	frame=single,
	captionpos=b,
	breaklines=true,
	breakatwhitespace=false,
	tabsize=4,
	showstringspaces=false,
	morekeywords={uint8_t, uint16_t, uint32_t, uint64_t, int8_t, int16_t, int32_t, int64_t},
}

% Set up the default behavior of listings
\lstset{
	basicstyle=\ttfamily,
	keywordstyle=\ttfamily,
	identifierstyle=\ttfamily,
	commentstyle=\ttfamily,
	stringstyle=\ttfamily,
	captionpos=b,
	breaklines=true,
	breakatwhitespace=false,
	tabsize=4,
	showstringspaces=false
}

% New commands
\newcommand{\abs}[1]{\left|#1\right|}
\newcommand{\textoverline}[1]{$\mkern 1.5mu\overline{\mkern-1.5mu\mbox{#1}\mkern-1.5mu}\mkern 1.5mu$}
%\newcommand{\truetype}[1]{{\texttt{#1}}}
\newcommand{\bashinline}[1]{\lstinline[language=bash]{#1}}
\newcommand{\bashinlinehl}[1]{\colorbox{linuxhighlightcolor}{\lstinline[language=bash]{#1}}}
\newcommand{\forthinline}[1]{\lstinline{#1}}
\newcommand{\forthinlinehl}[1]{\colorbox{linuxhighlightcolor}{\lstinline{#1}}}
\newcommand{\cinline}[1]{\lstinline[language=C]{#1}}
\newcommand{\cinlinehl}[1]{\colorbox{linuxhighlightcolor}{\lstinline[language=C]{#1}}}
\newcommand{\asminline}[1]{\texttt{#1}}
\newcommand{\asminlinehl}[1]{\colorbox{linuxhighlightcolor}{\texttt{#1}}}
\newcommand{\pin}[1]{\texttt{#1}}
\newcommand{\peripheral}[1]{\texttt{#1}}
\newcommand{\register}[1]{\texttt{#1}}
\newcommand{\bitfield}[1]{\texttt{#1}}

% Import defines
\newcommand{\AsicName}{Myshkin}
\newcommand{\AsicNameForUserGuide}{Myshkin}
\newcommand{\RevisionDateFull}{February 17$^\textrm{th}$, 2026}
\newcommand{\ProgaddrIrq}{0x9000}
\newcommand{\VectorsStartAddress}{0x8000}
\newcommand{\RamProgramStartAddress}{0x8100}
\newcommand{\RamStartAddress}{0x8000}
\newcommand{\RamEndAddress}{0xFFFF}
\newcommand{\RamEndWordAddress}{0xFFFC}
\newcommand{\RamNumPages}{128}
\newcommand{\PackageSize}{7}
\newcommand{\PackagePinCount}{44}
\newcommand{\PackagePinPitch}{0.5}
\newcommand{\PackagePinWidth}{0.25}
\newcommand{\PackagePinDepth}{0.4}
\newcommand{\PackageUnits}{mm}
\newcommand{\SpiFlashProgramAddress}{0x0000}


% Import tikzit style
%\input{murray.tikzstyles}
